The arc-midpoint lemma states that the midpoint $S$ of arc $AC$ not containing $B$ satisfies $SA=SC=SI$. Indeed $\angle SAI=\angle SAC+\angle CAI=\dfrac{B}{2}+\dfrac{A}{2}=90^\circ-\dfrac{C}{2}$, while $\angle AIS=180^\circ-\angle AIB=90^\circ-\dfrac{C}{2}$, so $SA=SI$; equal arcs give $SA=SC$. Similarly $TA=TI$. Thus $ST$ is the perpendicular bisector of $AI$. Because $AI$ bisects $\angle KAL$, reflection in $AI$ interchanges its side rays and fixes $ST$, so it interchanges $K,L$, giving $AK=AL$. Points on $ST$ are equidistant from $A,I$, so $KA=KI$ and $LA=LI$. All four sides of $AKIL$ are equal; hence it is a rhombus.
\begin{asy}
import graph;
usepackage("amsmath");
import geometry;
size(7cm);
pair A=(0,sqrt(3)),B=(-1,0),C=(1,0),I=(0,sqrt(3)/3),K=(-1/3,2sqrt(3)/3),L=(1/3,2sqrt(3)/3),S=(1,2sqrt(3)/3),T=(-1,2sqrt(3)/3);draw(A--B--C--cycle);draw(circle(I,2sqrt(3)/3));draw(S--T);draw(A--K--I--L--cycle,linewidth(1));draw(B--S,dashed);draw(C--T,dashed);label("$A$",A,(0,1));label("$B$",B,(-1,-1));label("$C$",C,(1,-1));label("$I$",I,(0,-1));label("$K$",K,(-1,0));label("$L$",L,(1,0));label("$S$",S,(1,0));label("$T$",T,(-1,0));
\end{asy}

